\chapter{Literature review}

This section will provide a literature review on the field of AI for multiphysics \& multiscale modeling.

Simulations are fundamental across science and engineering, but they are often conducted in isolation. For example, climate models predicting coastal erosion rarely account for human impacts like urbanization or mining, and are typically limited to specific regions and timescales. However, natural systems consist of numerous physical processes that operate at diverse spatial and temporal scales. For simulations to be both accurate and meaningful, they must incorporate multi-physics and multiscale dynamics. Similarly, AI and machine learning (ML), although powerful tools for handling multi-modal, multi-fidelity data, are limited by their data-driven approach. ML often overlooks fundamental physical laws, potentially leading to ill-posed problems and non-physical solutions. To ensure that AI and ML applications in science are reliable and meaningful, they must integrate multi-scale, multi-physics data and reveal underlying mechanisms driving complex behaviors.

\textbf{Multiphysics modeling} enables simulation-based analysis of interacting physical forces, such as fluid, thermal, structural, and electromagnetic effects, which operate simultaneously in real-world systems. By combining multiple physics solvers within a unified computational framework, multiphysics models allow engineers to examine complex behaviors across entire systems—going beyond the limitations of single-physics models. This approach is crucial in applications like fluid-structure interactions in aircraft, thermal-optics in automotive displays, and electromagnetic-thermal coupling in motors. By providing a realistic, system-level perspective, multiphysics solutions improve design accuracy, identify potential issues early in development, and enhance product performance across sectors such as aerospace, automotive, healthcare, and industry. While computationally intensive, advances in specialized software, like Ansys Fluent\textsuperscript{\textregistered} and System Coupling\texttrademark, have made it possible to manage complex simulations within a single platform, thereby simplifying the integration of high-fidelity, cross-domain models. As sustainable design and high-power-density electronics become increasingly important, multiphysics modeling is poised to play an even greater role in product development.

\textbf{Multiscale phenomena}, present in both natural systems and engineered models, involve interactions across multiple spatial and temporal scales. Examples include the hierarchical organization of society, the decomposition of signals into various frequencies (e.g., Fourier or wavelet expansions), and the microscale structure of materials influencing macroscale properties. In physics and engineering, multiscale problems arise in wave propagation, composite materials, and fluid dynamics, where small-scale structures or high Reynolds numbers create complex behaviors like boundary layers and turbulence. Multiscale models are essential in fields like biology and energy, connecting molecular to ecosystem levels and nanoscale materials to macroscopic processes. Although powerful, multiscale modeling is computationally demanding and often complicated by unknown physics, inhomogeneities, and complex interfaces.

We will discuss the following subfields of AI for multiphysics \& multiscale modeling: physics-informed neural network, physics-infused deep learning, and applications to science and industry.

{\Large \textbf{Physics-informed machine learning}} 
Physics-informed machine learning is characterized by imposing observational biases, inductive biases and learning biases to the neural network so that it can steer the neural network towards physically consistent solutions.

Here observational biases mainly refer to the training data. Imposing observational biases involves feeding training data to the neural network so that it can learn the underlying patterns and physical laws of the data. Inductive biases usually refer to model-based structure such as a differential equation that can be integrated with the neural network. Learning biases means choosing appropriate loss functions, constraints and inference algorithms to force the training of an ML system to converge to a physically consistent solution. 

The most representative architecture of physics-informed machine learning is the physics-informed neural network (PINN) \cite{raissi2019physics}. PINNs uniquely combine observational data with physical laws, making them highly effective for addressing complex problems where data may be incomplete, sparse, or noisy. By integrating the governing equations directly into the neural network’s loss function, PINNs are capable of producing solutions that respect physical constraints, even when conventional numerical methods struggle. This quality proves invaluable in applications involving ill-posed problems, where traditional techniques fail due to missing initial or boundary conditions. Another significant advantage of PINNs is their mesh-free nature, which allows them to bypass the complexities of generating computational grids, especially in irregular or dynamically changing domains. Their flexibility further extends to high-dimensional problems, such as those found in climate modeling \cite{lutjens2021pce}, biological systems \cite{arzani2021uncovering}, and material science \cite{manav2024phase}, where they can model complex dependencies among multiple parameters.

Despite these advantages, PINNs face certain limitations. Training neural networks to solve physics-based problems is computationally intensive, and handling high-dimensional PDEs exacerbates this demand \cite{karniadakis2021physics}. Additionally, while PINNs excel in many applications, they can encounter difficulties with problems involving high-frequency or oscillatory solutions, which require specific adaptations, such as Fourier feature networks, to achieve accuracy \cite{karniadakis2021physics}. Another challenge lies in balancing the physics-based loss terms to ensure convergence, a process that is often problem-dependent and can necessitate fine-tuning \cite{karniadakis2021physics}. This difficulty can make the training process more intricate and computationally expensive, especially for larger-scale or more complex problems.

The applications of PINNs span diverse scientific domains. In engineering, PINNs have been applied to reconstruct fluid dynamics in cases where direct measurements are limited, such as in medical imaging for blood flow estimation \cite{kissas2020machine} or in environmental modeling for air quality assessment \cite{hettige2024airphynet}. They are also valuable in physics and materials science, where they enable inverse modeling to infer unknown properties of materials from indirect measurements. PINNs are particularly well-suited for multiscale, multiphysics problems, making them instrumental in fields like geophysics \cite{zhu2021general}, climate prediction \cite{verma2024climode}, and fusion energy research \cite{mcdevitt2024physics}. With their scalable, flexible framework, PINNs are poised to play a key role in the future of scientific computing, although advances in algorithmic efficiency and tailored architectures will be necessary to overcome current limitations and maximize their potential in even more challenging applications.

{\Large \textbf{Physics-infused machine learning}}
Physics-infused machine learning could be viewed as a bidirectional version of Physics-informed machine learning. The latter emphasizes imposing physics constraint to improve the learning efficiency of neural networks. The former involves not only incorporating physical information but also utilizing neural networks to learn the underlying physics of a system. Examples of this line of work include \textit{Universal differential equations \cite{rackauckas2020universal}, Neural differential equations \cite{kidger2022on}, Kolmogorov-Arnold neural networks \cite{liu2024kan}, neural operators \cite{li2021fourier}, etc.} We will introduce each of them in this section.

\textbf{Universal differential equation} 
Universal Differential Equations (UDEs) are emerging as a powerful approach within scientific machine learning (SciML), merging traditional physical models with data-driven neural networks to improve flexibility and accuracy in scientific modeling. UDEs incorporate universal approximators, such as neural networks, into differential equations, either partially or fully, to enhance modeling capabilities where classical methods may fall short.

Mathematically, UDEs extend differential equations by embedding universal approximators that can capture unknown dynamics or adapt the model to better fit data. In its generalized form, a UDE can be expressed as 
\[
N[u(t), u(\alpha(t)), W(t), U_{\theta}(u, \beta(t))] = 0,
\]
where \(U_{\theta}\) represents the universal approximator (e.g., neural networks) that learns from data and captures complex relationships. This formulation allows UDEs to handle scenarios involving stochasticity, delays, and other complex behaviors that may be challenging for traditional models.

Applications of UDEs span a variety of scientific fields. They are used in biological systems to uncover hidden mechanisms \cite{rackauckas2020universal}, and in engineering for control systems. For instance, UDEs have been employed to solve high-dimensional Hamilton-Jacobi-Bellman equations \cite{rackauckas2020universal}, facilitating computationally efficient control in high-dimensional spaces. They also support symbolic regression \cite{rackauckas2020universal}, enabling researchers to discover or refine governing equations from data, which is valuable for understanding systems where the full physical model is unknown.

UDEs offer notable advantages over traditional methods. By integrating data-driven neural networks, UDEs capture complex dynamics that are difficult to model analytically, improving model accuracy and adaptability. They also reduce data requirements in cases where purely data-driven models would typically require large datasets, leveraging physical knowledge as an inductive bias. Additionally, UDEs are compatible with advanced numerical techniques, such as adaptive solvers and sensitivity analysis, making them computationally efficient and stable even in complex simulations.

Despite their advantages, UDEs have limitations. The integration of neural networks can lead to increased computational cost, especially when training deep architectures within UDEs. Furthermore，neuronal networks in UDEs might suffer from stability difficulties, mainly in stiff equations. UDEs require careful tuning of neural network architectures and training algorithms to ensure that the models are not only accurate but also computationally feasible. Lastly, interpreting neural networks within UDEs remains challenging, as the learned dynamics are often complex and not immediately transparent.

In conclusion, UDEs represent a significant advancement in scientific modeling, combining the strengths of traditional physics-based approaches with the flexibility of machine learning. This hybrid approach is especially powerful for tackling complex, multiscale problems in science and engineering, though the computational and interpretative challenges require further research to fully harness their potential.

\textbf{Neural differential equation}
The conjoining of differential equations and deep learning has long been considered as an important direction for development of AI. The broad field of Neural differential equation can be classified into many subfields such as Neural ODE \cite{chen2018neural}, Neural SDE \cite{li2020scalable}, Neural CDE \cite{kidger2020neural} and Neural rough differential equation \cite{salimans2022neural}, etc. They are newly arising yet powerful architectures for "AI for science". Based on their efficiency and impact, I consider them to be on par with famous AI architecture such as transformer \cite{vaswani2017attention}, convolutional neural network \cite{lecun1989backpropagation} and PINN. Neural ordinary differential equations (Neural ODEs), neural stochastic differential equations (Neural SDEs), and neural controlled differential equations (Neural CDEs) have recently gained attention as advanced tools for dynamic system modeling within machine learning. By combining neural networks with differential equation frameworks, these models provide flexibility in capturing complex dynamics, including those involving randomness or irregular time-series data. The following introduces each model, covering its formulation, applications, and strengths and limitations.

Neural ODEs model a system's evolution over time using ordinary differential equations (ODEs) integrated with neural networks. Their formulation is expressed as \( \frac{dy(t)}{dt} = f_\theta(t, y(t)) \), where \( f_\theta \) is a neural network parameterized by \(\theta\). In this setup, hidden states evolve based on a continuous-time function, allowing the calculation of hidden states at any time \(t\) by integrating the differential equation from an initial state \(y(0) = y_0\). For applications such as time-series data modeling, image classification and generative modeling, neural ODEs are very helpful. They offer a natural framework for continuous-time transformations and are commonly applied in time-series forecasting and continuous normalizing flows \cite{grathwohl2018ffjord}, where data undergoes continuous transformations. A primary advantage of Neural ODEs is their efficient memory usage since intermediate states can be reconstructed as needed, and they also support variable-step solvers that enhance computational efficiency. However, solving the ODE numerically can lead to increased training times, especially when dealing with stiff equations.

Neural SDEs extend Neural ODEs by incorporating randomness, making them suitable for modeling systems influenced by stochastic noises. Their formulation is typically represented by \( dy(t) = \mu(t, y(t)) \, dt + \sigma(t, y(t)) \circ dw(t) \), where \( \mu \) represents the drift (deterministic part), \( \sigma \) represents the diffusion (stochastic part), and \( w(t) \) denotes Brownian motion. The noise term \( \sigma \circ dw \) introduces stochasticity, making Neural SDEs well-suited for systems with inherent randomness. Applications of Neural SDEs are found in finance \cite{wiese2020quant} and population dynamics \cite{mckendrick2022neural} where uncertainty is an integral aspect of the system's evolution. In finance, for example, they model asset price dynamics, capturing both trend and volatility. Neural SDEs also aid in uncertainty estimation, as the stochastic component provides a distribution over possible outcomes rather than a single trajectory, making them beneficial for probabilistic forecasting. While Neural SDEs allow for capturing system uncertainty, their training can be more complex. They require specialized solvers and often entail greater computational costs due to the stochastic nature of the model, which can complicate gradient estimation during optimization.

Neural CDEs \cite{kidger2020neural} address path-dependent or sequential data, such as time-series data with irregular sampling. They use controlled differential equations (CDEs) in which the evolution of the system depends on both the hidden state and an external control signal \(x\): \( dy(t) = f_\theta(y(t)) \, dx(t) \). Here, \( f_\theta \) is the neural network vector field, and \( x(t) \) acts as the control path, enabling Neural CDEs to excel in sequence modeling. Applications of Neural CDEs include financial data analysis \cite{kidger2020neural} and healthcare, where observations may be irregularly sampled. They work effectively with sequential data that changes over time, adapting dynamically to irregularly observed data. In this way, Neural CDEs provide a continuous-time alternative to recurrent neural networks, allowing more natural representation of processes that evolve continuously. Neural CDEs are advantageous due to their adaptability to irregularly sampled data and their resemblance to recurrent networks in a continuous-time format. However, similar to Neural SDEs, they can be computationally intensive and require careful parameterization to maintain stability in the training process.

In summary, Neural ODEs, Neural SDEs, and Neural CDEs are valuable tools for modeling dynamic systems, each tailored to different data types and application scenarios. Neural ODEs excel in deterministic systems and continuous data transformations, Neural SDEs are suited for probabilistic modeling of systems with inherent randomness, and Neural CDEs handle sequence-dependent data with irregular sampling. These models continue to evolve as differential equation-solving techniques advance, opening new applications and further improving their robustness and efficiency.

\textbf{Kolmogorov-Arnold neural network}
The Kolmogorov-Arnold Neural Network (KAN) builds on a foundational concept in mathematics, specifically the Kolmogorov-Arnold superposition theorem, which states that any multivariate continuous function can be represented as a finite sum of compositions of univariate functions. In recent developments, this theorem has inspired new approaches to solving partial differential equations (PDEs) through machine learning, specifically within the framework of Physics-Informed Machine Learning (PIML). This has led to the development of Physics-Informed Kolmogorov-Arnold Networks (PIKANs) \cite{wang2024kolmogorov}, which are highly specialized for handling complex PDEs across various scientific and engineering domains.

In Kolmogorov-Arnold Networks, the function representation is structured through a two-layer composition: an outer layer composed of univariate functions and an inner layer consisting of multivariate inputs. The network takes input in the form of multivariate data \( z^{(l-1)} \) from the previous layer \( l-1 \) and transforms it using both univariate and multivariate functions. Specifically, each output node of the layer \( z^{(l)} \) is calculated as:
\[
z^{(l)} = \sum_{i=1}^{H_{2}} \Phi_i \left( \sum_{j=1}^{H_{1}} \varphi_{i,j}(z^{(l-1)}_j) \right)
\]
where \( \Phi_i \) and \( \varphi_{i,j} \) represent the outer and inner univariate functions, respectively. This structure allows KANs to approximate high-dimensional functions efficiently, especially in the context of solving differential equations with complex boundary conditions and variable coefficients.

Kolmogorov-Arnold Networks have shown significant promise in physics-informed learning, where they can be adapted to solve both forward and inverse problems related to PDEs. In forward problems, PIKANs are used to approximate the solution of PDEs by encoding physical laws directly into the network architecture. For instance, they can be configured to model heat transfer, fluid dynamics, and other complex systems where classical methods struggle due to high dimensionality or nonlinearity. Additionally, the networks can also solve inverse problems, where unknown parameters or hidden physical fields are inferred from partial data. For example, applications have included inferring material properties in elasticity problems and reconstructing flow fields in fluid dynamics.

PIKANs offer several advantages over traditional PIML architectures. Firstly, they provide more compact representations than standard neural networks, requiring fewer parameters to achieve similar levels of accuracy. This efficiency arises from their two-level structure, which reduces the burden of capturing high-dimensional relationships in complex systems. Moreover, PIKANs are inherently modular, making them suitable for decomposing large domains into smaller, manageable sub-domains. This domain decomposition capability facilitates parallel processing and is particularly useful for handling extensive computations, such as those encountered in 3D fluid simulations. Another notable advantage is that KANs offer improved numerical stability when approximating solutions to PDEs. Unlike traditional deep networks, KANs are less prone to instabilities related to vanishing or exploding gradients, which can otherwise impede training and hinder convergence. This robustness makes PIKANs particularly useful in cases where the solutions exhibit sharp gradients or other complexities that challenge regular neural network architectures.

Although Kolmogorov-Arnold Networks have their strengths, they do have some limitations. One of the primary challenges is the selection of suitable inner and outer functions \( \Phi_i \) and \( \varphi_{i,j} \), which often requires significant domain knowledge and experimentation. Training PIKANs can also be resource-intensive, particularly when dealing with high-dimensional challenges that involve complex boundary conditions. Although KANs mitigate some issues related to scalability, their implementation still demands substantial computational resources. Furthermore, while PIKANs are robust in representing continuous functions, they may not be as effective when dealing with highly discontinuous solutions. Discontinuities require specialized handling, such as adaptive refinement or hybrid methods, to achieve accurate results. Lastly, the interpretability of KANs, though improved over traditional deep learning models, can still be challenging in practice, as the precise relationship between the network parameters and the modeled physical phenomena is not always straightforward.

In summary, Kolmogorov-Arnold Networks and their physics-informed variant, PIKANs, represent a powerful extension of PIML for solving PDEs. Their architecture, inspired by mathematical theory, is well-suited to capturing high-dimensional relationships while preserving computational efficiency. The integration of PIKANs into scientific machine learning opens new avenues for addressing challenges in physics-driven simulations, offering both practical advantages and areas for further refinement.

{\Large \textbf{Applications of AI to physics}}

\textbf{Deep Flame: combustion simulation platform \cite{mao2023deepflame}}
Combustion, especially in multiphase and turbulent environments, involves the complex integration of multiscale problems and has long been a challenging area for large-scale scientific computation. Computational fluid dynamics (CFD), based on the numerical solution of the macroscopic continuous Navier-Stokes equations, has become irreplaceable in fields like aerospace and weather forecasting. Numerical discretization is necessary to solve continuous partial differential equations computationally. The finite volume method based on spatial grid discretization, one of the most widely used approaches, has been implemented and industrialized in various commercial CFD software abroad, such as Fluent and StarCCM+.

On the other hand, macroscopic-level chemical reaction kinetics calculations commonly use the Arrhenius formula to determine elementary reaction rates. Combustion systems typically involve hundreds of components and even more elementary reactions, requiring the solution of a large system of ordinary differential equations over time. For solving such systems, there are similar commercial software solutions abroad, such as ChemkinPro \cite{designs2011chemkin} and CosiLab.

However, fluid dynamics simulations involving combustion reactions require coupling both CFD and chemical reaction kinetics, presenting even greater challenges. Although various commercial software have made significant strides in functionality, current levels of computational accuracy and efficiency remain insufficient for industrial applications.

Under the AI for Science research paradigm, the DeepFlame project organizes open-source platforms like OpenFOAM \cite{jasak2007openfoam}, Cantera, and Torch \cite{collobert2002torch}, and leverages next-generation computational infrastructure such as heterogeneous parallelism and AI accelerators. The aim is to develop a high-precision, high-efficiency, easy-to-use, and widely applicable numerical simulation program for reactive flows in combustion. This project, initiated and led by Chen Zhi from Peking University, with AI modeling support from Zhang Tianhan at Southern University of Science and Technology and Xu Zhiqin at Shanghai Jiao Tong University, is supported by the Beijing Academy of Artificial Intelligence (AISI) as part of their collaborative research platform and the DeepModeling community for code hosting and related technical support.

\textbf{FourCastNet: large-scale weather prediction with the assistance of neural operator \cite{pathak2022fourcastnet}}
FourCastNet, short for Fourier Forecasting Neural Network, is a global, data-driven weather forecasting model designed for high-resolution predictions using the Adaptive Fourier Neural Operator (AFNO) \cite{guibas2021adaptive}. Developed by a team including researchers from NVIDIA and several academic institutions, FourCastNet achieves a resolution of 0.25° (approximately 30 km) and has proven effective for short- to medium-term forecasts of atmospheric variables such as surface wind speed, precipitation, and atmospheric water vapor. This model has shown particular strengths in applications that benefit from fine-scale atmospheric details, including disaster response and wind energy planning. FourCastNet’s performance matches the European Centre for Medium-Range Weather Forecasts (ECMWF) Integrated Forecasting System (IFS) for larger-scale forecasts, while surpassing IFS on small-scale phenomena like precipitation, due to the model's ability to resolve finer details.

Traditional numerical weather prediction (NWP) models require intensive computation, often demanding extensive computing resources and time. In contrast, FourCastNet leverages deep learning to achieve significantly faster forecast times. For instance, it can produce a week-long forecast in less than two seconds, representing orders of magnitude faster processing than traditional models. This rapid forecasting enables large ensemble simulations, where multiple forecast scenarios can be generated affordably and quickly to support probabilistic forecasting, crucial for predicting and preparing for extreme weather events such as tropical cyclones.

FourCastNet integrates a Fourier-based token-mixing scheme with a Vision Transformer (ViT) \cite{dosovitskiy2020image} backbone. This architecture enables high-resolution predictions by effectively handling long-range spatial dependencies, making it well-suited for representing complex atmospheric structures. The ViT component is known for efficiently modeling interactions across large distances in data, which is critical in high-resolution weather models that must capture both global and local weather dynamics. The Adaptive Fourier Neural Operator adds a layer of scalability and resolution independence, improving FourCastNet’s ability to generalize across varying scales.

The FourCastNet model was trained using the ECMWF ERA5 reanalysis dataset, a comprehensive global weather dataset. Key variables in the training set included wind speeds, temperatures, surface pressures, and humidity levels at various atmospheric levels, chosen to capture the critical interactions that affect weather patterns. While FourCastNet’s training focused on these fundamental atmospheric properties, its architecture also allows it to model complex processes like precipitation through a separate diagnostic model. This approach simplifies the main model while preserving the accuracy needed for reliable precipitation forecasting, a particularly challenging variable due to its sporadic and localized nature.

By providing both high-resolution and rapid forecasts, FourCastNet opens up new possibilities in meteorology, especially in probabilistic forecasting. Massive ensemble forecasts, which are essential for capturing uncertainties and refining extreme weather predictions, can be generated using FourCastNet within seconds, unlike the computationally expensive ensemble forecasting traditionally employed by NWP models. Additionally, the energy efficiency of FourCastNet is notable, as it consumes significantly less energy per forecast than traditional models, which could prove beneficial for sustainable computing in weather prediction.


% ------------------------------------------------------------------------


%%% Local Variables: 
%%% mode: latex
%%% TeX-master: "../thesis"
%%% End: 
